% ---------------------------------------------------------------------------
% Conceptual diagrams.
%
% Presentation only. Nothing in this file carries a number, an estimate, or an
% interval. It defines the drawing language for the thesis conceptual figures so
% that a hand-drawn TikZ diagram and a matplotlib artifact can sit on the same
% page without a clash of colour, weight, or type.
%
% The palette is the LaTeX twin of src/viz/style.py. Keep the two in step.
% ---------------------------------------------------------------------------

\usepackage{tikz}
\usetikzlibrary{positioning,arrows.meta,calc,fit,backgrounds,shapes.geometric}

% ------------------------------------------------------------------ palette
% Mirrors src/viz/style.py: navy carries the estimate, steel and its fill carry
% uncertainty, brick is the single contrast accent, the greys carry structure.
\definecolor{thesisnavy}{HTML}{1F3B57}
\definecolor{thesisnavylight}{HTML}{41627F}
\definecolor{thesissteel}{HTML}{AFC2D1}
\definecolor{thesisfill}{HTML}{DCE5EC}
\definecolor{thesisbrick}{HTML}{9C4722}
\definecolor{thesisgrey}{HTML}{6E7B85}
\definecolor{thesisgreylight}{HTML}{AAB4BB}
\definecolor{thesisrule}{HTML}{3F484E}
\definecolor{thesisink}{HTML}{22282C}

% ------------------------------------------------------------------- styles
% Diagram type is the document serif, one step down from the body, with the 1.5
% body leading relaxed so a wrapped label reads as a compact block rather than as
% loose running text. The body sets a low hyphenation penalty, which is right for
% a justified 16 cm measure and wrong inside a 30 mm box, so short labels break
% between words instead of splitting them.
\newcommand{\thesisdiagramtype}{%
  \setstretch{1.06}\footnotesize
  \hyphenpenalty=10000
  \exhyphenpenalty=10000
  \pretolerance=9999
  \tolerance=9999
  \emergencystretch=0pt
}

\tikzset{
  % Picture-level defaults.
  thesisdiagram/.style={
    line width=0.5pt,
    draw=thesisrule,
    text=thesisink,
    font=\thesisdiagramtype,
    node distance=6mm,
    inner sep=0pt,
    outer sep=0pt,
  },
  % Soft box: hairline rule, barely tinted fill, generous interior padding.
  softnode/.style={
    draw=thesisrule,
    line width=0.5pt,
    fill=thesisfill!45!white,
    rounded corners=2pt,
    inner xsep=6.5pt,
    inner ysep=6pt,
    outer sep=0pt,
    align=center,
    text=thesisink,
    font=\thesisdiagramtype,
  },
  % Faint variant. Used where a stage is read alongside an argument rather than
  % asserted as part of it.
  softnode faint/.style={
    softnode,
    draw=thesisgreylight,
    dash pattern=on 1.7pt off 1.5pt,
    fill=white,
    text=thesisrule,
  },
  % Fixed-width chain box, so every stage in a chain is the same size.
  chainnode/.style={
    softnode,
    text width=#1,
    minimum height=17mm,
  },
  chainnode faint/.style={
    softnode faint,
    text width=#1,
    minimum height=17mm,
  },
  % Connectors. Thin, short, and the same grey as the box rules.
  chainarrow/.style={
    draw=thesisrule,
    line width=0.5pt,
    -{Stealth[length=4.2pt,width=3.0pt,round]},
    shorten <=0.6pt,
    shorten >=0.6pt,
  },
  % Faint connector for a link that is interpretive rather than causal.
  chainarrow soft/.style={
    chainarrow,
    draw=thesisgrey,
    dash pattern=on 1.7pt off 1.5pt,
  },
}

% -------------------------------------------------------- reusable chain
% \thesismechanismchain[<stage text width>]{stage 1}{stage 2}{stage 3}{stage 4}
%
% Four equally sized stages, evenly spaced, joined by short arrows. The terminal
% stage takes the faint treatment because it is evidence read alongside the
% chain, not a strict downstream consequence of it.
\newcommand{\thesismechanismchain}[5][0.183\linewidth]{%
  \begin{tikzpicture}[thesisdiagram]
    \node[chainnode={#1}]                        (stageone)   {#2};
    \node[chainnode={#1}, right=of stageone]     (stagetwo)   {#3};
    \node[chainnode={#1}, right=of stagetwo]     (stagethree) {#4};
    \node[chainnode faint={#1}, right=of stagethree] (stagefour) {#5};
    \draw[chainarrow]      (stageone)   -- (stagetwo);
    \draw[chainarrow]      (stagetwo)   -- (stagethree);
    \draw[chainarrow soft] (stagethree) -- (stagefour);
  \end{tikzpicture}%
}
